\section{Inequalities, representations, and separation}\label{sec:preliminaries}

\subsection{Cut inequalities}

Let $V=\{0\}\cup W$ be a finite vertex set with a distinguished root $0$.
Write $E(V)$ for the unordered pairs of distinct vertices.  A cut vector
$d^S\in\{0,1\}^{E(V)}$, for $S\subseteq V$, has $d^S_{ij}=1$ exactly when
one of $i,j$ belongs to $S$.  The cut polytope is
$\CUT(V)=\conv\{d^S:S\subseteq V\}$.
For $d\in\R^{E(V)}$ and $b\in\Z^V$, put
\begin{equation}\label{eq:cut-forms}
 \sig(b)=\sum_{i\in V}b_i,\qquad
 Q(b,d)=\sum_{\{i,j\}\in E(V)}b_i b_jd_{ij},\qquad
 \gam(b)=\min_{s\in\{\pm1\}^V}\abs{s\trans b}.
\end{equation}
We call $\gam(b)$ the gap of $b$.  It has the same parity as $\sig(b)$,
and $0\leq\gam(b)\leq\abs{\sig(b)}$.

\Cref{tab:classes} fixes the names used in this paper.  The positive
semidefinite (psd) inequalities allow real coefficients; the other rows use
integer coefficients.  These conventions follow the comparison of the
families in \citet{GKL2012}.  In particular, a rounded positive
semidefinite inequality need not be hypermetric: its coefficient sum can
be any odd integer.  The $k$-gonal class in the table fixes the
$\ell_1$ norm of the coefficient vector, rather than its sum.
We call $\sum_i\abs{b_i}$ its gonality.

\begin{table}[t]
\centering\small
\caption{Cut inequality families.  The last row defines one class for each
positive integer $k$.}\label{tab:classes}
\begin{tabular}{@{}lll@{}}
\toprule
Family & Coefficients & Inequality\\
\midrule
Hypermetric & $b\in\Z^V$, $\sig(b)=1$ & $Q(b,d)\leq0$\\
Negative type & $b\in\Z^V$, $\sig(b)=0$ & $Q(b,d)\leq0$\\
Positive semidefinite & $b\in\R^V$ & $Q(b,d)\leq\sig(b)^2/4$\\
Rounded positive semidefinite & $b\in\Z^V$, $\sig(b)$ odd
 & $Q(b,d)\leq(\sig(b)^2-1)/4$\\
Gap & $b\in\Z^V$ & $Q(b,d)\leq(\sig(b)^2-\gam(b)^2)/4$\\
Gap-$1$, gap-$0$ & $\gam(b)=1$, $\gam(b)=0$ & As in the gap row\\
Odd clique & $b\in\{0,\pm1\}^V$, $\sig(b)$ odd
 & $Q(b,d)\leq(\sig(b)^2-1)/4$\\
Pure hypermetric & $b\in\{0,\pm1\}^V$, $\sig(b)=1$ & $Q(b,d)\leq0$\\
$k$-gonal & $b\in\Z^V$, $\sum_i\abs{b_i}=k$
 & $Q(b,d)\leq\lfloor\sig(b)^2/4\rfloor$\\
\bottomrule
\end{tabular}
\end{table}

Every inequality in the table is valid for $\CUT(V)$.  Indeed, on a cut
$S$ let $a=\sum_{i\in S}b_i$.  Then
$Q(b,d^S)=a(\sig(b)-a)=(\sig(b)^2-(2a-\sig(b))^2)/4$.
The gap bound follows by minimizing the absolute signed sum.
The other bounds follow from this identity, parity, or the nonnegativity
of a square.  The same argument for real $b$ proves the positive
semidefinite bound.  Hypermetric coefficients have gap $1$, so
\begin{equation}\label{eq:containment}
 \text{hypermetric}\ \subseteq\ \text{gap-$1$}\ \subseteq\
 \text{rounded positive semidefinite}.
\end{equation}
Containment alone does not transfer hardness of separation: a larger
family may cut both the yes and the no instances of a reduction.

Let $D$ be the symmetric matrix with zero diagonal and off-diagonal
entries $d_{ij}$, and let $J=\one\one\trans$.  Set
\begin{equation}\label{eq:Z}
 Z(d)=J-2D.
\end{equation}
Then $b\trans Z(d)b=\sig(b)^2-4Q(b,d)$.  The elliptope
$\ELL(V)=\{Z\succeq0:\diag Z=\one\}$ therefore enforces every positive
semidefinite inequality.  Its relative interior consists of its positive
definite matrices.  We also use the metric polytope $\MET(V)$, defined
by the triangle inequalities $d_{ij}\leq d_{ik}+d_{jk}$ and perimeter
inequalities $d_{ij}+d_{ik}+d_{jk}\leq2$ for distinct $i,j,k$.
For $\abs V\geq3$ these imply $0\leq d_{ij}\leq1$.

\subsection{Correlation and moment representations}

The following linear bijection sends a distance vector to a symmetric
matrix on $W$:
\begin{equation}\label{eq:M}
 M(d)_{ii}=d_{0i},\qquad
 M(d)_{ij}=\frac{d_{0i}+d_{0j}-d_{ij}}2\quad(i\ne j).
\end{equation}
Its inverse is $d_{0i}=M_{ii}$ and
$d_{ij}=M_{ii}+M_{jj}-2M_{ij}$.  Write $\delta=\diag M$ and
\begin{equation}\label{eq:g}
 g_M(z)=z\trans Mz-\delta\trans z.
\end{equation}
For every $b=(b_0,z)\in\R^V$, direct substitution gives
\begin{equation}\label{eq:Q-M}
 Q(b,d)=\sig(b)\delta\trans z-z\trans Mz.
\end{equation}
In particular, $z\mapsto(1-\one\trans z,z)$ is a bijection from
$\Z^W$ to the hypermetric coefficient vectors, and
\begin{equation}\label{eq:hypermetric-g}
 Q((1-\one\trans z,z),d)=-g_M(z).
\end{equation}
Thus hypermetric separation asks whether an integer $z$ satisfies
$g_M(z)<0$.  The inequalities $g_M(z)\geq0$, $z\in\Z^W$, are called
hypermetric correlation inequalities.  Negative-type membership is equivalent to $M\succeq0$,
so a violation exists exactly when $M$ has a negative direction:
the vectors with coefficient sum zero give $Q(b,d)=-z\trans Mz$.
Although the negative-type row of \cref{tab:classes} uses integers,
a rational matrix with a negative real direction also has a rational,
and hence an integer, negative direction.

If $M$ is the Gram matrix of vectors $u_i$ and there is a point $c$ with
$2\ip{c}{u_i}=\norm{u_i}^2$ for each $i$, then
\begin{equation}\label{eq:sphere}
 g_M(z)=\norm{\sum_i z_i u_i-c}^2-\norm c^2.
\end{equation}
The points $0,u_i$ lie on the sphere with center $c$ through $0$.
A violated hypermetric inequality is an integer combination strictly
inside this sphere.  For rational $M$, these combinations form a
lattice in their span, as follows from \cref{lem:rational-quotient}.
For positive definite $M$, the center in the span of the
$u_i$ always exists.  This geometric description underlies the earlier
lattice approach to hypermetricity \citep{AvisGrishukhin1993}; our
reduction specifies the sphere and all of its interior lattice points.

For the Boolean quadric representation, use variables
$x_i$ and $y_{ij}$, $i,j\in W$, $i<j$, and put
\begin{equation}\label{eq:Y}
 Y(x,B)=\begin{pmatrix}1&x\trans\\x&B\end{pmatrix},\qquad
 B_{ii}=x_i,\quad B_{ij}=y_{ij}\quad(i\ne j).
\end{equation}
The Boolean quadric polytope $\BQP(W)$ is the convex hull of
$(x,y)=(z,(z_i z_j)_{i<j})$ for $z\in\{0,1\}^W$.
The covariance map
\begin{equation}\label{eq:covariance}
 d_{0i}=x_i,\qquad d_{ij}=x_i+x_j-2y_{ij}
\end{equation}
is an affine bijection between $\BQP(W)$ and $\CUT(V)$
\citep{Padberg1989,Letchford2022}.  For arbitrary points in these
ambient spaces, $B=M(d)$ and $x=\diag B$.
Moreover,
\begin{equation}\label{eq:congruence}
 Z(d)=L Y(x,B)L\trans,\qquad
 L=\begin{pmatrix}1&0\\\one&-2I\end{pmatrix}.
\end{equation}
Consequently $Y\succ0$ if and only if $Z(d)\succ0$.

Some algorithms below apply more generally to a rational moment matrix
$Y$ with $Y_{00}=1$, without the identities $B_{ii}=x_i$.
We call this the general integer moment setting.  For $v\in\Z^{\{0\}\cup W}$
define the split slack
\begin{equation}\label{eq:q}
 q_Y(v)=v\trans Yv+e_0\trans Yv.
\end{equation}
It is the linearization of
$(v_0+v_W\trans z)(v_0+v_W\trans z+1)$, which is nonnegative for every
integer $z$.  In Boolean data these are precisely the Boros--Hammer
inequalities \citep{BorosHammer1993,LetchfordSorensen2012}.
For $w\in\Z^W$ and $s\in\Z$, it is convenient to write their slack as
\begin{equation}\label{eq:BH}
 h_Y(w,s)=w\trans Bw-(2s+1)x\trans w+s(s+1)
          =q_Y((s,-w))=q_Y((-s-1,w)).
\end{equation}
The two data pairs $(w,s)$ and $(-w,-s-1)$ define the same polynomial.
With
\begin{equation}\label{eq:beta}
 \beta(w,s)=(2s+1-\one\trans w,w)
\end{equation}
and Boolean data, \cref{eq:Q-M} gives
\begin{equation}\label{eq:BH-rounded}
 h_Y(w,s)=\frac{\sig(\beta(w,s))^2-1}{4}-Q(\beta(w,s),d).
\end{equation}
This is a bijective correspondence between Boros--Hammer data and
rounded positive semidefinite coefficients.  Equivalently, for
$u=e_0+2v$ one has
\begin{equation}\label{eq:parity-split}
 q_Y(v)=\frac{u\trans Yu-1}{4}.
\end{equation}

For later use, the signed-coefficient Boros--Hammer family takes
disjoint subsets $S,T\subseteq W$ and $w=\one_S-\one_T$.
Dividing \cref{eq:BH} by $2$ yields
\begin{equation}\label{eq:finite-BH}
 s\sum_{i\in S}x_i+\sum_{i\in S,j\in T}y_{ij}
 \ \leq\ (s+1)\sum_{i\in T}x_i
 +\sum_{\{i,j\}\subseteq S}y_{ij}
 +\sum_{\{i,j\}\subseteq T}y_{ij}+\frac{s(s+1)}2.
\end{equation}
We call the $T=\varnothing$, $0\leq s\leq\abs S-1$ subfamily the
clique inequalities, and the $s=0$ subfamily the cut inequalities.
These are the established Boolean quadric families surveyed by
\citet{Letchford2022}.  Their violations in this divided form are half
the corresponding cut-form violations in \cref{eq:BH-rounded}.

\subsection{Switching and the decision problem}

Switching a distance vector on $U\subseteq V$ replaces $d_{ij}$ by
$1-d_{ij}$ when exactly one endpoint lies in $U$, and leaves it
unchanged otherwise.  Write this vector as $d^U$ and set
$\varepsilon_i=-1$ for $i\in U$, $\varepsilon_i=1$ otherwise.
Then $Z(d^U)=\diag(\varepsilon)Z(d)\diag(\varepsilon)$.
Switching preserves the cut and metric polytopes and positive
definiteness of $Z$.  It sends a rounded positive semidefinite
inequality with vector $b$ to one with vector
$\diag(\varepsilon)b$: parity is unchanged.
It also preserves the gap and the restrictions $b_i\in\{0,\pm1\}$.
It does not in general preserve the condition $\sig(b)=1$ defining
the hypermetric class.

All complexity statements concern rational input, encoded in binary.
Let $L$ denote its total bit length.  Exact separation is the decision
problem of whether at least one inequality in the specified family is
strictly violated.  A separation algorithm must also return such an
inequality on a yes instance.  There is no numerical tolerance in this
definition.  The complementary membership problem asks whether every
inequality of the family is satisfied.

For a rational-input problem, strong NP-hardness here means hardness
already on reductions in which every numerator and denominator is
bounded by a polynomial in the combinatorial size of the source
instance.  The bounds below are stronger: each reduction point has a
common denominator bounded by such a polynomial.  For infinite
families, membership in $\NP$ also requires a violated inequality with
polynomial bit length; \cref{sec:algorithms} proves the required bounds.
Promise statements such as $d\in\MET(V)$ and $Z(d)\succ0$ specify
properties of all constructed inputs, independently of whether they
are yes or no instances.

A running time $f(r)\poly(L)$ is fixed-parameter tractable in $r$:
the exponent of $L$ is independent of $r$.  An XP bound permits that
exponent to depend on the parameter and therefore gives polynomial
time for each fixed parameter value.  We distinguish these two
guarantees throughout.
